\section{Formulation}
\label{sec:formulation}

\status{Ready to write: the equations and the non-dimensionalisation convention
are established and checked against the code. The text below is complete on the
points that matter for reproducibility.}

\subsection{Governing equations}

Incompressible two-phase flow with surface tension and a polymeric
contribution. In non-dimensional form,
\begin{equation}
\rho\left(\frac{\partial \mathbf{u}}{\partial t}
          + \mathbf{u}\cdot\nabla\mathbf{u}\right)
= -\nabla p
  + \frac{1}{\mathrm{Re}}\nabla\cdot\!\left[\mu\left(\nabla\mathbf{u}
    + \nabla\mathbf{u}^{\mathsf{T}}\right)\right]
  + \rho\,\hat{\mathbf{g}}
  + \frac{1}{\mathrm{Bo}}\,\kappa\mathbf{n}\delta
  + \nabla\cdot\mathbf{S},
\label{eq:momentum}
\end{equation}
with
\begin{equation}
\mathrm{Re} = \frac{\rho_1 g^{1/2} L^{3/2}}{\mu_1},
\qquad
\mathrm{Bo} = \frac{\rho_1 g L^2}{\sigma},
\qquad
V = \sqrt{gL},
\qquad
L = 2r,
\label{eq:dimensionless}
\end{equation}
where $L$ is the bubble \emph{diameter} and the characteristic velocity is tied
to gravity. The volume fraction $f$ obeys
\begin{equation}
\frac{\partial f}{\partial t} + \nabla\cdot(\mathbf{u}f) = 0,
\label{eq:transport}
\end{equation}
and the properties are interpolated as $\rho = (1-f)\rho_0 + f\rho_1$ and
$\mu = (1-f)\mu_0 + f\mu_1$.

\subsection{The polymeric tensor is not the stored tensor}

This deserves its own subsection because it is a source of silent error. The
solver stores $\mathbf{S}$; the polymeric tensor is
\begin{equation}
\mathbf{T} = \mathbf{S} + \frac{2(1-\beta)}{\mathrm{Re}}\,\mathbf{D},
\qquad
\mathbf{D} = \tfrac{1}{2}\!\left(\nabla\mathbf{u}
           + \nabla\mathbf{u}^{\mathsf{T}}\right),
\label{eq:T-from-S}
\end{equation}
and $\beta$ is the \emph{interpolated} value between the phases, by the same
expression the solver uses:
\begin{equation}
\beta = \frac{(1-f)\,\mu_0\beta_0 + f\,\mu_1\beta_1}{\mu}.
\label{eq:beta-interp}
\end{equation}
Mistaking $\mathbf{S}$ for $\mathbf{T}$ raises no error; it produces a
plausible and wrong field. The output files of this work carry both, labelled.

\subsection{Constitutive model}

Oldroyd-B, as an evolution equation for the conformation tensor $\mathbf{A}$:
\begin{equation}
\mathrm{De}\left(\frac{\partial \mathbf{A}}{\partial t}
  + \mathbf{u}\cdot\nabla\mathbf{A}
  - \nabla\mathbf{u}\cdot\mathbf{A}
  - \mathbf{A}\cdot\nabla\mathbf{u}^{\mathsf{T}}\right)
= \mathbf{I} - \mathbf{A},
\qquad
\mathrm{De} = \frac{U\lambda}{L}.
\label{eq:oldroyd}
\end{equation}
Oldroyd-B has no finite extensibility limit, so the stress grows without bound
under sustained extension. That is a property of the model, not of the
discretisation, and it bounds how far the viscoelastic cases can be pushed ---
Section~\ref{sec:viscoelastic} returns to it.

\subsection{The solver's non-dimensionalisation convention}
\label{sec:convention}

This subsection exists for reproducibility, and its content was \emph{read from
the code and checked numerically}, not assumed. The solver integrates the
equations in \textbf{dimensional} variables, with the dimensionless groups set
so that the coefficients of (\ref{eq:momentum}) reduce to the physical values:
\begin{equation}
\mathrm{Re} = 1,
\qquad
\mathrm{Fr} = 1/\sqrt{g},
\qquad
\mathrm{Ca} = 1/\sigma,
\end{equation}
with $\rho$ and $\mu$ entering at their tabulated values. The gravitational
term appears as $1/\mathrm{Fr}^2$ and the interfacial one as
$1/(\mathrm{Ca}\cdot\mathrm{Re})$. Checked: $\mathrm{Fr} = 1.010153$ gives
$1/\mathrm{Fr}^2 = 0.979999$, which is $g = 0.98$; and
$\mathrm{Ca} = 0.0408163265$ gives $1/(\mathrm{Ca}\cdot\mathrm{Re}) = 24.5$,
which is $\sigma$. With $U = L = 1$ at this scaling, $\mathrm{De} = \lambda$.
